\BourbakiExerciseGroup

\begin{enumerate}
\def\labelenumi{\arabic{enumi})}
\tightlist
\item
  If \(x\) and \(y\) are real, we have
\end{enumerate}

\[
[ x + y ] = [ x ] + [ y ] + \varepsilon \quad \text { where } \varepsilon = 0 \quad \text { or } \quad \varepsilon = 1,
\]

\[
[ x - y ] = [ x ] - [ y ] - \varepsilon \quad \text { where } \varepsilon = 0 \quad \text { or } \quad \varepsilon = 1,
\]

\[
[ x ] + [ x + y ] + [ y ] \leqslant [ 2 x ] + [ 2 y ],
\]

\[
[ x ] + \left[ x + \frac {1}{n} \right] + \left[ x + \frac {2}{n} \right] + \dots + \left[ x + \frac {n - 1}{n} \right] = [ n x ],
\]

\[
\left[ \frac {[ n x ]}{n} \right] = [ x ]
\]

for any integer \(n > 0\) .

\begin{enumerate}
\def\labelenumi{\arabic{enumi})}
\setcounter{enumi}{1}
\item
  Let \((\varepsilon_{n})\) be a strictly decreasing sequence of finite numbers \(>0\) which tend to 0. For each \(x \in \mathbb{R}\) , let \(r_{n}(x)\) be the multiple of \(\varepsilon_{n}\) which is the approximation by defect to \(x\) to within \(\varepsilon_{n}\) . Then the sequence \((r_{n}(x))\) is increasing for all \(x \in \mathbb{R}\) if and only if \(\varepsilon_{n}\) is an integer multiple of \(\varepsilon_{n+1}\) for each \(n\) .
\item
  Let \(a\) be an integer \(> 1\) . A real number \(x\) is rational if and only if its expansion to base \(a\) , say \(x = p_0 + \sum_{n=0}^{\infty} u_n a^{-n}\) , is periodic, that is to say if and only if there exist two integers \(n_0\) and \(r > 0\) such that \(u_{n+r} = u_n\) for all \(n \geqslant n_0\) .
\item
  Let \((u_{n})\) be a summable sequence of numbers \(>0\) in \(\mathbf{R}\) which satisfies the following conditions: \(u_{n+1} \leqslant u_{n}\) and \(u_{n} \leqslant \sum_{k=1}^{\infty} u_{n+k}\) for all \(n \geqslant 0\) . Show that for every number \(a\) such that \(0 < a \leqslant s = \sum_{n=0}^{\infty} u_{n}\) , there exists a subset I of \(\mathbf{N}\) such that \(a = \sum_{n \in \mathbf{I}} u_{n}\) . These conditions are satisfied in particular if \(u_{n+1} \leqslant u_{n} \leqslant 2u_{n+1}\) for all \(n\) . Consider the case of dyadic expansions.
\item
  For each real number \(x \in ]0, I]\) , there is a unique infinite increasing sequence \((q_n)\) of integers \(> 0\) such that \(x = \sum_{n=1}^{\infty} I / (q_1 q_2 \ldots q_n)\) . The number \(x\) is rational if and only if \(q_{n+1} = q_n\) from a certain value of \(n\) onwards.
\item
  For each real number \(x > \mathbf{I}\) , there is a unique infinite sequence \((t_n)\) of integers \(\geqslant \mathbf{I}\) such that \(t_{n+1} \geqslant t_n^2\) for all \(n \geqslant \mathbf{I}\) and such that \(x = \prod_{n=1}^{\infty} \left( \mathbf{I} + \frac{\mathbf{I}}{t_n} \right)\) (use Exercise 19 of §7). The number \(x\) is rational if and only if \(t_{n+1} = t_n^2\) form a certain value of \(n\) onwards. {[}To show that the condition is necessary, show that if we put \(x_n = \prod_{k=n+1}^{\infty} \left( \mathbf{I} + \frac{\mathbf{I}}{t_k} \right)\) then the denominators of the fractions \(\frac{x_n}{x_n - \mathbf{I}}\) (in their irreducible form) form a decreasing sequence{]}.
\end{enumerate}

¶ * 7) Let \((\varepsilon_n)_{n \geqslant 0}\) be a sequence of numbers each of which is 1 or --- 1. Show that the real number

\[
x _ {n} = \varepsilon_ {0} \sqrt {2 + \varepsilon_ {1} \sqrt {2 + \varepsilon_ {2} \sqrt {2 + \cdots + \varepsilon_ {n} \sqrt {2}}}}
\]

is defined and is equal to

\[
2 \sin \left(\frac {\pi}{4} \sum_ {k = 0} ^ {n} \frac {\varepsilon_ {0} \varepsilon_ {1} \dots \varepsilon_ {k}}{2 ^ {k}}\right).
\]

Deduce that \(\lim_{n\to \infty}x_n\) exists and that for each real number \(x\) , such that \(-2\leqslant x\leqslant 2\) , there is a sequence \((\varepsilon_n)\) such that \(x\) is equal to the limit of the corresponding sequence \((x_{n})\) . For what values of \(x\) is the sequence \((\varepsilon_n)\) unique? For what values of \(x\) is it periodic? (see Exercise 3).*

\begin{enumerate}
\def\labelenumi{\arabic{enumi})}
\setcounter{enumi}{7}
\item
  Show that there is no non-constant mapping \(\psi\) of an interval \(\mathbf{I} \subset \mathbf{R}\) into the set \(\mathbf{N}^{\mathbb{N}}\) of all sequences of natural numbers which has the following property for each \(x \in \mathbf{I}\) : for each integer \(n \geqslant 0\) there is a neighbourhood \(\mathbf{V}\) of \(x\) in \(\mathbf{I}\) such that, for each \(y \in \mathbf{V}\) , the first \(n\) terms of the sequence \(\psi(y)\) are the same as the first \(n\) terms of the sequence \(\psi(x)\) (remark that this would imply the continuity of \(\psi\) , if each factor \(\mathbf{N}\) carries the discrete topology and \(\mathbf{N}^{\mathbb{N}}\) the product topology).
\item
  Show that Cantor's triadic set K (§ 2, no. 5) is the set of all real numbers \(x \in [0, 1]\) whose triadic expansion (resp. improper triadic expansion if \(x\) is the origin of an interval contiguous to K)
\end{enumerate}

\[
x = \sum_ {n = 1} ^ {\infty} u _ {n} 3 ^ {- n}
\]

is such that \(u_{n} \neq \mathrm{I}\) for each \(n\) (so that \(u_{n} = 0\) or \(u_{n} = 2\) ). If \(v_{n} = \frac{\mathrm{I}}{2} u_{n}\) for each \(n\) , and if \(f(x) = \sum_{n=1}^{\infty} v_{n} 2^{-n}\) , show that \(f\) is a surjective continuous mapping of \(\mathbf{K}\) onto \([\mathrm{o}, \mathrm{I}]\) ; deduce that \(\mathbf{K}\) has the power of the continuum.

¶ 10) Let \((d_n)\) be a base sequence, let \(a_n = d_n / d_{n-1}\) ( \(n \geqslant 1\) ), let \((n_i)\) be a strictly increasing sequence of integers, and let \((b_i)\) be a sequence of integers such that \(0 < b_i < a_{n_i} - 1\) for each \(i\) . Let \(G\) be the set of all \(x \in [0, 1]\) such that, if \(\sum_{n=1}^{\infty} u_n / d_n\) is the expansion of \(x\) , or the improper expansion if it exists, then \(u_{n_i} \neq b_i\) for all \(i \in N\) . Show that \(G\) is a totally disconnected perfect set, and that if \(l\) is the sum of the lengths of the intervals contiguous to \(G\) and contained in \([0, 1]\) , then

\[
\mathrm{I} - l = \prod_ {i = 0} ^ {\infty} \left(\mathrm{I} - \frac {\mathrm{I}}{a _ {n _ {i}}}\right).
\]

¶ 11) Let X be a Hausdorff topological space. Suppose that, for each finite sequence s whose terms are equal to o or 1, there is a nonempty subset \(\mathrm{A}(s)\) of X which satisfies the following conditions: (i) If s is a sequence of n terms (each of which is o or i) and \(s'\) , \(s''\) are the sequences of \(n + 1\) terms (each of which is o or i) whose first n terms are the same as those of s, then \(\mathrm{A}(s) = \mathrm{A}(s') \cup \mathrm{A}(s'')\) ; and if \(s_0\) is the empty sequence, \(\mathrm{A}(s_0) = \mathrm{X}\) .

\begin{enumerate}
\def\labelenumi{(\roman{enumi})}
\setcounter{enumi}{1}
\tightlist
\item
  For each infinite sequence \((u_n)_{n\geqslant 0}\) whose terms are equal to 0 or 1, if \(s_n\) denotes the finite sequence \((u_k)_{0\leqslant k\leqslant n}\) , then the filter base formed by the \(\mathbf{A}(s_n)\) converges to a point of \(\mathbf{X}\) .
\end{enumerate}

Show that under these conditions there is a surjective continuous mapping of Cantor's triadic set K onto X.

¶ 12) Deduce from Exercise 11 that:

\begin{enumerate}
\def\labelenumi{\alph{enumi})}
\item
  If A is a compact subset of R, there is a continuous mapping of Cantor's triadic set K onto A {[}take the A(s) to be the intersections of A with suitably chosen intervals{]}.
\item
  If also A is perfect and totally disconnected, it is homeomorphic to K {[}same method, by arranging that if \(s_{1}\) and \(s_{2}\) are two distinct sequences with the same number of terms (equal to o or 1) then
\end{enumerate}

\[
\mathbf {A} (s _ {1}) \cap \mathbf {A} (s _ {2}) = \varnothing_ {1} ].
\]

¶ 13) Let X be a countable Hausdorff space. Suppose that, for each finite sequence s all of whose terms are o or 1, there is a non-empty subset B(s) of X which satisfies the following conditions:

\begin{enumerate}
\def\labelenumi{(\roman{enumi})}
\item
  If \(s\) is a finite sequence of \(n\) terms (equal to \(\mathbf{o}\) or \(\mathbf{i}\) ) and \(s', s''\) are the sequences of \(n + \mathbf{i}\) terms (equal to \(\mathbf{o}\) or \(\mathbf{i}\) ) whose first \(n\) terms are the same as those of \(s\) , then \(\mathbf{B}(s') \cup \mathbf{B}(s'') = \mathbf{B}(s)\) and \(\mathbf{B}(s') \cap \mathbf{B}(s'') = \emptyset\) ; and if \(s_0\) is the empty sequence, \(\mathbf{B}(s_0) = \mathbf{X}\) .
\item
  For each \(x \in \mathbf{X}\) , if \(s_n\) is the (unique) sequence of \(n\) terms equal to 0 or 1 such that \(x \in \mathrm{B}(s_n)\) , then the filter base formed by the \(\mathrm{B}(s_n)\) converges to \(x\) in \(\mathbf{X}\) .
\end{enumerate}

Show that, under these conditions, E is homeomorphic to the rational line Q. {[}Show first that X is homeomorphic to a countable subset of Cantor's triadic set K which is dense in K and contains no end-point of an interval contiguous to K. In order to do this, remark that the hypotheses associate with each \(x \in X\) an infinite sequence \((u_{n}(x))\) whose terms are equal to 0 or 1; using the fact that X is countable, show that the bijection \(s \to \mathbf{B}(s)\) can be modified in such a way that for no \(x \in E\) do the \(u_{n}(x)\) form a stationary sequence; finally use Exercise 10 of §2{]}.

Deduce that every countable subspace of \(\mathbf{R}\) which has no isolated points is homeomorphic to \(\mathbf{Q}\) .

\begin{enumerate}
\def\labelenumi{\arabic{enumi})}
\setcounter{enumi}{13}
\tightlist
\item
  Show that every closed subset of \(\mathbf{R}\) is either countable or has the power of the continuum {[}use Exercise 12 b) above and Exercise 16 of Chapter I, § 9{]}.
\end{enumerate}


\begin{enumerate}
\def\labelenumi{\arabic{enumi})}
\setcounter{enumi}{14}
\item
  Show that the set of open subsets of \(\mathbf{R}\) and the set of compact, totally disconnected, perfect subsets of \(\mathbf{R}\) have the power of the continuum.
\item
  \begin{enumerate}
  \def\labelenumii{\alph{enumii})}
  \tightlist
  \item
    Let A be a countable closed subset of R and let f be a continuous real-valued function defined on R. Show that if f is constant in each of the intervals contiguous to A, then f is constant (use Bolzano's theorem).
  \end{enumerate}
\end{enumerate}

\begin{enumerate}
\def\labelenumi{\alph{enumi})}
\setcounter{enumi}{1}
\tightlist
\item
  If \(f\) is the continuous mapping of Cantor's triadic set \(\mathbf{K}\) onto \([0, 1]\) defined in Exercise 9, show that \(f\) can be extended to a continuous function on \(\mathbf{R}\) which is constant on each of the intervals contiguous to \(\mathbf{K}\) .
\end{enumerate}

\begin{enumerate}
\def\labelenumi{\arabic{enumi})}
\setcounter{enumi}{16}
\tightlist
\item
  Let X be a topological space which has a countable dense subset. Show that the set of all real-valued continuous functions on X has the power of the continuum.
\end{enumerate}
% semantic-fragments: legacy-input-seam
\relax{}
